% Proved necessary profile after complete root and companion audits; no realization or exclusion claim.
\subsection{The necessary ramification profile of the surviving full image}

Keep the hypotheses and normalization of B87:
\(p=5\), \(e_K=20\), \(\mu_5\subset K\),
algebraically closed residue field, common tame good
field \(L/K\) of degree 24, and a different-special-character
actual \(125\)-comparison with tame reducible residual inertia.
The parameters are \(B_1=25,B_2=75\), with high lines
\(\mathbf F_5e_1,\mathbf F_5e_2\), respectively.
B87 gives \(W_1=M_2(\mathbf F_5)\). Let
\(F_2=L(M[25])\) be the actual common second division field,
so \(\operatorname{Gal}(F_2/L)=W_1\).
All breaks below use normalized valuations and the convention
\(b(g)=v(g\varpi-\varpi)-1\).

\begin{proposition}[Both actual point-field ramification filtrations]
\label{prop:B-RD88-point-fields}
\status{Proved; full independent root and companion audits passed.}
Let \(E_i=L(P_i)\), where \(P_i\) is any primitive
second-level point with \(5P_i=e_i\).
Each \(E_i/L\) is Galois of degree 25.
Its two lower breaks and its two coordinate cyclic
quotient breaks are necessarily
\[
\begin{array}{c|c|c|c}
 &\text{lower breaks}&\text{diagonal quotient}&\text{off-diagonal quotient}\\\hline
 E_1/L&(6,336)&72&6\\
 E_2/L&(24,234)&24&66.
\end{array}                                             \tag{B-RD88.1}
\]
The diagonal quotients here are different actual fields;
they are the \(A_{11}\) and \(A_{22}\) quotients of \(W_1\).
\end{proposition}
\begin{proof}
Evaluation on \(e_i\) is surjective. Since the first-level
module is \(L\)-rational, its kernel is the actual stabilizer
of every such lift. It is a normal subgroup of the additive
second-level image, proving the field assertion. In \(E_1\)
the translation coordinates are \((A_{11},A_{21})\);
in \(E_2\) they are \((A_{12},A_{22})\).
Their tame weights are respectively \((0,6)\) and
\((18,0)\). The tame congruence and norm-quotient arguments
of Lemma~\ref{lem:B-RD87-fields} apply separately to each
field, without identifying their diagonal quotients.

Write \(d_i\) and \(o_i\) for their diagonal and
off-diagonal cyclic quotient breaks. Thus
\[
 d_i\in24\mathbf Z_{>0},\qquad
 o_1\equiv6\pmod{24},\qquad o_2\equiv18\pmod{24}.
                                                        \tag{B-RD88.2}
\]
There are two distinct lower breaks \(t<u\). The last
ramification line is an eigenaxis, and the larger cyclic
quotient break is \((u+4t)/5\), while the smaller is \(t\).
The four actual coordinate rows are still (B-RD87.7),
with the same strict cutoff \(12000\).

We use explicitly the fixed-coordinate minimum from
(B-RD87.3). For a row with leading valuation \(A\),
let \(J\) be its first nonzero exponent prime to five.
At lower break \(b\) the first possible displacement
degree is \(\min(A+5b,J+b)\); a strict smaller minimum
than the actual displacement can be removed only by
the tie \(J=A+4b\). Every prime exponent present
satisfies \(J>A\), and \(5\nmid J\).

First consider \(E_1\). An off-diagonal last line
would make the first row of (B-RD87.7) give
\(o_1=6>d_1\), impossible by (B-RD88.2).
The last line is therefore diagonal, and the second
row gives \(d_1=72\). Consequently
\[
       t=o_1\in\{6,30,54\},\qquad u=360-4t.          \tag{B-RD88.3}
\]
For the second row, whose high-displacement line is not
the last line, the compulsory prime degree is
\(J=15+4t\). If \(t=30\), it is 135, divisible by
five, a contradiction. If \(t=54\), apply the first
row instead: its displacement outside the last line
is 125, whereas \(A+5t=95+270=365\).
It would require \(J=125-54=71<A=95\), impossible.
Only \(t=6,u=336\) remains.

Now consider \(E_2\). If the last line were diagonal,
the third row would give \(d_2=24\). Necessarily
\(o_2=t=18\) and \(u=120-72=48\).
In the fourth row the smaller-break elements have
displacement 375, while \(A+5t=45+90=135\).
Cancellation forces \(J=45+4\cdot18=117\).
At the last line's break 48, this prime term changes
uniquely first, at degree \(117+48=165\), less than
the other possible degree \(45+5\cdot48=285\).
Its actual displacement would have to be 1125,
a contradiction.

Thus the last line in \(E_2\) is off-diagonal.
The fourth row gives \(o_2=66\), and hence
\[
             t=d_2\in\{24,48\},\qquad u=330-4t.      \tag{B-RD88.4}
\]
If \(t=48,u=138\), the third row has displacement
125 outside its high line, whereas \(A+5t=245\).
It follows that \(J=125-48=77\). At the last line,
the unique first change would then be
\(77+138=215<5+5\cdot138=695\), contradicting its
actual displacement 2375. Therefore \(t=24,u=234\).
This proves every entry of (B-RD88.1). All first-change
comparisons were strictly below \(v_{E_i}(5)=12000\);
none assumes an unproved independence of higher coefficients.
\end{proof}

\begin{corollary}[The full second division field and its cyclotomic quotient]
\label{cor:B-RD88-cyclotomic}
\status{Proved; full independent root and companion audits passed.}
The upper breaks of \(F_2/L\), together with the
corresponding coordinate characters, are
\[
\begin{array}{c|rrrr}
 \text{character}&A_{21}&A_{22}&A_{12}&A_{11}\\\hline
 \text{upper break}&6&24&66&72.
\end{array}                                             \tag{B-RD88.5}
\]
Its lower breaks are \(6,96,1146,1896\), with successive
ramification ranks \(4,3,2,1\) over \(\mathbf F_5\).
Among its 156 cyclic degree-five subfields, the numbers
with breaks \(6,24,66,72\) are \(1,5,25,125\), respectively.
The actual cyclotomic quotient has
\[
 b\bigl(L(\mu_{25})/L\bigr)=72,
 \qquad b\bigl(K(\mu_{25})/K\bigr)=3.                \tag{B-RD88.6}
\]
Both displayed extensions have degree five.
\end{corollary}
\begin{proof}
The four coordinate cyclic quotients have exactly
the breaks in (B-RD88.5), by the proposition.
We recall why a nonzero sum of independent cyclic
characters with different breaks has the larger break.
Factor through the actual degree-25 quotient defined
by the two characters. Its two quotient breaks differ,
so it has two ramification breaks and a last line.
The smaller-break character vanishes on the last line,
whereas the larger-break one does not. Their sum also
does not, and the same exact coset-average formula
gives it the larger break. This proves the assertion
directly, without an equality of unrelated point fields.
Induction now shows that every nonzero combination of
the four coordinate characters has break equal to the
largest break of its nonzero coordinates.

The quotient compatibility of upper numbering therefore
gives (B-RD88.5) and the stated successive ranks. Here
the complete norm-polynomial/Herbrand argument is the
one already checked in
\cite[Proposition 3.7 and Theorem 3.10, full proofs,
pp.~14--16]{ColmezCorpsLocaux}; the normalization is
also that of Serre's original Section 3.1.
Integrating the indices \(1,5,25,125\) gives the lower
breaks explicitly:
\[
 6,\quad6+5(24-6)=96,\quad
 96+25(66-24)=1146,\quad
 1146+125(72-66)=1896.                                 \tag{B-RD88.7}
\]
The number of one-dimensional character spaces whose
last nonzero coordinate is the \(j\)th in that order
is \(5^{j-1}\). This gives all 156 cyclic subfields.

On \(L\), the determinant of the actual good height-two
representation is the cyclotomic character. Modulo 25,
\(\det(1+5A)=1+5\operatorname{tr}(A)\).
Thus \(L(\mu_{25})\) is precisely the quotient field
for \(A_{11}+A_{22}\), which has break 72 by the
unequal-break sum argument. It is nontrivial, and
\(\mu_5\subset K\), so it and \(K(\mu_{25})/K\)
have degree five.

For clarity, the final tame scaling is an actual
field calculation. Put \(C=K(\mu_{25})\) and
\(D=LC\). The extensions \(C/K\) and \(L/K\)
are Galois of coprime degrees 5 and 24. Hence
\(D/C\) is tame of degree 24 and
\(\operatorname{Gal}(D/K)\) is their direct product.
Choose a uniformizer \(x\) of \(C\). The tame
uniformizer/Hensel argument provides a uniformizer
\(z\) of \(D\) with \(z^{24}=x\); see the complete
argument of \cite[Proposition 2.12, p.~11]{ColmezCorpsLocaux}.
For a nonidentity \(g\in\operatorname{Gal}(D/L)\),
the ratio \(g(z)/z\) has residue one, and
\[
   (g(z)/z)^{24}=g(x)/x.
\]
As 24 is a unit, the valuations of these ratios minus
one agree. Thus
\(b(D/L)=24b(C/K)\). The left side is 72, giving
the asserted break three over \(K\).
\end{proof}

\begin{remark}[A necessary profile, not a final exclusion]
\label{rem:B-RD88-scope}
\status{Proved; full independent root and companion audits passed.}
This result constrains the actual remaining \((4,4)\)
image in the tame reducible sector at \(p=5,e_K=20\).
It neither constructs an ambiguous pair nor proves that
the required cyclotomic break three is impossible over
the allowed base fields. In particular, the fixed value
\(24^2\) from B86 concerned the different hypothesis
\(\mu_{25}\subset K\) and is not reused here.
The two diagonal quotient fields are distinct because
\(A_{11},A_{22}\) are independent characters on the
actual full first layer. No false shared-diagonal-field
assertion is imported from the three-dimensional case.

The only original-source inputs beyond the earlier
proved dependencies are the complete Colmez tame
uniformizer and ramification quotient/Herbrand proofs,
already freshly read and scoped in B86. No additional
finite-residue assumption or priority claim is made.

The exact retained supplement is
\path{code/reduction_descent/verify_p5_full_image_ramification.py}.
Its first author run passed at 01:46:08.611352 UTC
on 20 November 2026 in 0.001338374990154989 seconds.
It checks both initially excluded last-line cases,
all three remaining \(E_1\) residue choices and both
\(E_2\) choices, the four upper-to-lower intervals,
all 624 nonzero coordinate characters and their
156 projective classes, the trace quotient and tame
scaling, and the equality of the different computed
from the lower filtration and from the character
conductors. These are exact finite consequences of
the preceding proof; no bounded search for an
unknown ramification break replaces that proof.
No failed mathematical execution preceded this run;
its exact code and output remain retained.

\smallskip
\noindent\emph{Independent audits.}
The complete root and companion proof, source and code audits passed.
Their literal author-output-preserving reruns took
\(0.0014191469963407144\) and
\(0.0013393660046858713\) seconds, respectively.
The exact audited candidate, full read scopes, code,
outputs and promotion diff are retained.
\end{remark}
